
\newcommand{\nomedoc}{Modello}
\newcommand{\versione}{0.3}
\newcommand{\nomefile}{modello\versione.pdf}
\newcommand{\datacreazione}{2 Dicembre 2010}
\newcommand{\datamodifica}{3 Dicembre 2010}
\newcommand{\stato}{formale}
\newcommand{\uso}{interno}
\newcommand{\redazione}{Mandolo Andrea}
\newcommand{\verifica}{Baffo}
\newcommand{\approvazione}{Valter}
\newcommand{\distribuzione}{
VT.G \\
& Prof. Vardanega Tullio }

% FUNZIONI TIPOGRAFICHE
\newcommand{\co}{\texttt} % courier
\newcommand{\bo}{\textbf} % bold
\newcommand{\pr}{\par\medskip} % paragrafo spaziato
\newcommand{\sca}{\textsc} % small caps

\documentclass[a4paper,12pt]{report}
% 10pt,11pt,12pt
% titlepage, notitlepage -> per dare inizio o no ad una nuova pagina dopo titolo
% twoside -> per dire se fronte-retro
\usepackage[latin1]{inputenc}
% per caratteri accentati
\usepackage[italian]{babel}
% per regole sintattiche italiane
\usepackage[bookmarks=true, pdfborder={0 0 0 0}]{hyperref}
% per collegamenti ipertestuali
\usepackage{graphicx}
% per inserimento immagini

% \usepackage{enumerate}
% per personalizzare elenchi puntati

\usepackage[hmargin=2cm]{geometry} %margine 2 cm
%\geometry{options varie}

% comandi per gestire meglio header e footer
\usepackage{fancyhdr}  % header e footer
\usepackage{totpages}
\pagestyle{fancy}
\renewcommand{\headrulewidth}{0.4pt}
\renewcommand{\footrulewidth}{0.4pt}

\setlength{\headheight}{1.2cm} % NON TOCCARE
\setlength{\voffset}{-1.5cm} % NON TOCCARE
\setlength{\textheight}{666pt} % NON TOCCARE
\setlength{\footskip}{60pt}
\setlength{\parindent}{0pt} % INDENTAZIONE

\lhead{\nomedoc\  (ver. \versione)}
\chead{}
\rhead{\includegraphics[height=1cm]{img/netmus.png}}
\lfoot{\includegraphics[height=0.8cm]{img/logo.png}}
\cfoot{}
\rfoot{\thepage}

\usepackage{titlesec}
\titleformat{\chapter}{\normalfont\huge\bfseries}
{\thechapter}{20pt}{\Huge}

\usepackage{rotating}   % PER TABELLE E AMBIENTI RUOTATI
\usepackage{array}
\usepackage{color}
\usepackage{colortbl}  % VARIE PER GESTIRE I COLORI
\definecolor{Orange}{RGB}{255,127,0}   % ARANCIO ACCES0
\definecolor{orange}{RGB}{255,207,80}  % ARANCIO TENUE

\addtocontents{toc}{\protect\thispagestyle{fancy}}  % PER INDICI CON + PAGINE
\usepackage[font=it]{caption}    % PER RENDERE CORSIVE LE DIDASCALIE
\usepackage{eurosym}  % PER SIMBOLO EURO

% \usepackage{listings}   per codice sorgente

\author{VT.G - Valter Texas Group}

\begin{document}

\pagenumbering{Roman} % INIZIO NUMERAZIONE ARABA

\vspace*{1cm}
\begin{center}

\begin{LARGE} \sca{Federico Baron} \end{LARGE}\\
\vspace{0.5cm}
\begin{Large}
\emph{fede.baron.89@gmail.com} \end{Large}\\
\vspace*{1cm} \includegraphics[width=5cm]{img/logo.png}\\
\vspace{0.5cm}
\begin{Large} \emph{``Comunicazione Aumentata/Alternativa per Giovani Ospiti
della Terapia Intensiva Pediatrica''} \end{Large}\\
\vspace{3cm}
\begin{Large} \sca{\nomedoc} \end{Large}\\
\end{center}
\vspace{1cm}

% INFORMAZIONI DOCUMENTO
\begin{center}
\begin{tabular}{r|l}
\hline & \\
\bo{Nome} & \nomefile \\
\bo{Versione attuale} & \versione \\
\bo{Data creazione} & \datacreazione \\
\bo{Data ultima modifica} & \datamodifica \\
\bo{Redazione} & \redazione \\
& \\\hline
\end{tabular}
\end{center}
\newpage

% REGISTRO MODIFICHE
\section*{Registro delle modifiche}
\begin{tabular}{lll}
% REVISIONI DEL DOCUMENTO
% DALLA PIU' NUOVA ALLA PIU' VECCHIA

% TEMPLATE PER REVISIONE
\bo{Data:} 03/12/2010 &
\bo{Versione:} 0.2 &
\bo{Autore:} Mandolo Andrea\\
\hline\\
\multicolumn{3}{p{470px}}{ Rinominato file, dimensione caratteri 12pt,
aggiunti alcuni contenuti comuni.}\\
\\

\bo{Data:} 03/12/2010 &
\bo{Versione:} 0.1 &
\bo{Autore:} Mandolo Andrea\\
\hline\\
\multicolumn{3}{p{470px}}{ Stesura prima versione di modello.}\\ \\

\end{tabular}

% INDICE
\tableofcontents
\thispagestyle{fancy} % per lo stile di header e footer


\chapter*{Sommario}


\thispagestyle{fancy} % serve perche' nelle pagine di inizio Chapter esca header e footer
\pagenumbering{arabic} % INIZIO NUMERAZIONE NORMALE
\rfoot{\thepage\ di \pageref{TotPages}}
\addcontentsline{toc}{chapter}{Sommario}

\chapter{Introduzione}
\thispagestyle{fancy} % serve perche' nelle pagine di inizio Chapter esca header e footer

\section{Scopo del documento}



\section{Scopo del prodotto}
Il progetto \underline{NetMus} nasce con lo scopo di realizzare un sistema
software basato su \underline{cloud} \underline{computing}, per memorizzare
informazioni di brani musicali in profili utente online.\\ Tali informazioni vengono estratte da
dispositivi musicali o di archiviazione \underline{USB} al momento della loro connessione.

\section{Glossario}
Il Glossario \`e definito con un documento a parte
(\emph{Glossario-\versioneglossario.pdf}). Tutti i termini caratterizzati da
\underline{questa sottolineatura} sono ivi definiti.\\
Verr\`a sottolineata solamente la prima occorrenza di ciascun
termine presente nel Glossario, per non compromettere la leggibilit\`a del documento.

\section{Riferimenti}

\subsection{Normativi} % oppure rif. a Norme di progetto con leggi e tutto
\begin{itemize}
  \item ISO/IEC 12207:1995 - Cicli di vita software
  \item ISO/IEC 9126:2001 - Quality Model
  \item \emph{NormeDiProgetto-\versionenormeprogetto.pdf} che regola e
  accompagna tutti i documenti ufficiali.
\end{itemize}
\newpage
\subsection{Informativi}
\begin{itemize}
  \item Capitolato d'appalto CO2-NETMUS del corso di Ingegneria del Software
  A.A. 2010/11 :\\
  \url{http://www.math.unipd.it/~tullio/IS-1/2010/Progetto/NetMus.pdf}
  \item Slide delle lezioni del corso:\\
  \url{http://www.math.unipd.it/~tullio/IS-1/2010/}
  \item Verbale intervista proponente:\\
  \co{allegato Verbale-1.0.pdf}
  \item Sistema di cloud Google App Engine:\\
  \url{http://code.google.com/intl/it/appengine/}
\end{itemize}


% INIZIO CAPITOLO 2

\end{document}